\documentclass[a4paper,12pt,twoside]{memoir}

% Castellano
\usepackage[spanish,es-tabla,english]{babel}
\selectlanguage{english}
\usepackage[utf8]{inputenc}
\usepackage[T1]{fontenc}
\usepackage{lmodern} % Scalable font
\usepackage{microtype}
\usepackage{placeins}

\RequirePackage{booktabs}
\RequirePackage[table]{xcolor}
\RequirePackage{xtab}
\RequirePackage{multirow}

% Links
\PassOptionsToPackage{hyphens}{url}\usepackage[colorlinks]{hyperref}
\hypersetup{
	allcolors = {red}
}

% Ecuaciones
\usepackage{amsmath}

% Rutas de fichero / paquete
\newcommand{\ruta}[1]{{\sffamily #1}}

% Párrafos
\nonzeroparskip

% Huérfanas y viudas
\widowpenalty100000
\clubpenalty100000

% Imágenes

% Comando para insertar una imagen en un lugar concreto.
% Los parámetros son:
% 1 --> Ruta absoluta/relativa de la figura
% 2 --> Texto a pie de figura
% 3 --> Tamaño en tanto por uno relativo al ancho de página
\usepackage{graphicx}
\newcommand{\imagen}[3]{
	\begin{figure}[!h]
		\centering
		\includegraphics[width=#3\textwidth]{#1}
		\caption{#2}\label{fig:#1}
	\end{figure}
	\FloatBarrier
}

% Comando para insertar una imagen sin posición.
% Los parámetros son:
% 1 --> Ruta absoluta/relativa de la figura
% 2 --> Texto a pie de figura
% 3 --> Tamaño en tanto por uno relativo al ancho de página
\newcommand{\imagenflotante}[3]{
	\begin{figure}
		\centering
		\includegraphics[width=#3\textwidth]{#1}
		\caption{#2}\label{fig:#1}
	\end{figure}
}

% El comando \figura nos permite insertar figuras comodamente, y utilizando
% siempre el mismo formato. Los parametros son:
% 1 --> Porcentaje del ancho de página que ocupará la figura (de 0 a 1)
% 2 --> Fichero de la imagen
% 3 --> Texto a pie de imagen
% 4 --> Etiqueta (label) para referencias
% 5 --> Opciones que queramos pasarle al \includegraphics
% 6 --> Opciones de posicionamiento a pasarle a \begin{figure}
\newcommand{\figuraConPosicion}[6]{%
  \setlength{\anchoFloat}{#1\textwidth}%
  \addtolength{\anchoFloat}{-4\fboxsep}%
  \setlength{\anchoFigura}{\anchoFloat}%
  \begin{figure}[#6]
    \begin{center}%
      \Ovalbox{%
        \begin{minipage}{\anchoFloat}%
          \begin{center}%
            \includegraphics[width=\anchoFigura,#5]{#2}%
            \caption{#3}%
            \label{#4}%
          \end{center}%
        \end{minipage}
      }%
    \end{center}%
  \end{figure}%
}

%
% Comando para incluir imágenes en formato apaisado (sin marco).
\newcommand{\figuraApaisadaSinMarco}[5]{%
  \begin{figure}%
    \begin{center}%
    \includegraphics[angle=90,height=#1\textheight,#5]{#2}%
    \caption{#3}%
    \label{#4}%
    \end{center}%
  \end{figure}%
}
% Para las tablas
\newcommand{\otoprule}{\midrule [\heavyrulewidth]}
%
% Nuevo comando para tablas pequeñas (menos de una página).
\newcommand{\tablaSmall}[5]{%
 \begin{table}
  \begin{center}
   \rowcolors {2}{gray!35}{}
   \begin{tabular}{#2}
    \toprule
    #4
    \otoprule
    #5
    \bottomrule
   \end{tabular}
   \caption{#1}
   \label{tabla:#3}
  \end{center}
 \end{table}
}

%
% Nuevo comando para tablas pequeñas (menos de una página).
\newcommand{\tablaSmallSinColores}[5]{%
 \begin{table}[H]
  \begin{center}
   \begin{tabular}{#2}
    \toprule
    #4
    \otoprule
    #5
    \bottomrule
   \end{tabular}
   \caption{#1}
   \label{tabla:#3}
  \end{center}
 \end{table}
}

\newcommand{\tablaApaisadaSmall}[5]{%
\begin{landscape}
  \begin{table}
   \begin{center}
    \rowcolors {2}{gray!35}{}
    \begin{tabular}{#2}
     \toprule
     #4
     \otoprule
     #5
     \bottomrule
    \end{tabular}
    \caption{#1}
    \label{tabla:#3}
   \end{center}
  \end{table}
\end{landscape}
}

%
% Nuevo comando para tablas grandes con cabecera y filas alternas coloreadas en gris.
\newcommand{\tabla}[6]{%
  \begin{center}
    \tablefirsthead{
      \toprule
      #5
      \otoprule
    }
    \tablehead{
      \multicolumn{#3}{l}{\small\sl continúa desde la página anterior}\\
      \toprule
      #5
      \otoprule
    }
    \tabletail{
      \hline
      \multicolumn{#3}{r}{\small\sl continúa en la página siguiente}\\
    }
    \tablelasttail{
      \hline
    }
    \bottomcaption{#1}
    \rowcolors {2}{gray!35}{}
    \begin{xtabular}{#2}
      #6
      \bottomrule
    \end{xtabular}
    \label{tabla:#4}
  \end{center}
}

%
% Nuevo comando para tablas grandes con cabecera.
\newcommand{\tablaSinColores}[6]{%
  \begin{center}
    \tablefirsthead{
      \toprule
      #5
      \otoprule
    }
    \tablehead{
      \multicolumn{#3}{l}{\small\sl continúa desde la página anterior}\\
      \toprule
      #5
      \otoprule
    }
    \tabletail{
      \hline
      \multicolumn{#3}{r}{\small\sl continúa en la página siguiente}\\
    }
    \tablelasttail{
      \hline
    }
    \bottomcaption{#1}
    \begin{xtabular}{#2}
      #6
      \bottomrule
    \end{xtabular}
    \label{tabla:#4}
  \end{center}
}

%
% Nuevo comando para tablas grandes sin cabecera.
\newcommand{\tablaSinCabecera}[5]{%
  \begin{center}
    \tablefirsthead{
      \toprule
    }
    \tablehead{
      \multicolumn{#3}{l}{\small\sl continúa desde la página anterior}\\
      \hline
    }
    \tabletail{
      \hline
      \multicolumn{#3}{r}{\small\sl continúa en la página siguiente}\\
    }
    \tablelasttail{
      \hline
    }
    \bottomcaption{#1}
  \begin{xtabular}{#2}
    #5
   \bottomrule
  \end{xtabular}
  \label{tabla:#4}
  \end{center}
}



\definecolor{cgoLight}{HTML}{EEEEEE}
\definecolor{cgoExtralight}{HTML}{FFFFFF}

%
% Nuevo comando para tablas grandes sin cabecera.
\newcommand{\tablaSinCabeceraConBandas}[5]{%
  \begin{center}
    \tablefirsthead{
      \toprule
    }
    \tablehead{
      \multicolumn{#3}{l}{\small\sl continúa desde la página anterior}\\
      \hline
    }
    \tabletail{
      \hline
      \multicolumn{#3}{r}{\small\sl continúa en la página siguiente}\\
    }
    \tablelasttail{
      \hline
    }
    \bottomcaption{#1}
    \rowcolors[]{1}{cgoExtralight}{cgoLight}

  \begin{xtabular}{#2}
    #5
   \bottomrule
  \end{xtabular}
  \label{tabla:#4}
  \end{center}
}



\graphicspath{ {./img/} }

% Capítulos
\chapterstyle{bianchi}
\newcommand{\capitulo}[2]{
	\setcounter{chapter}{#1}
	\setcounter{section}{0}
	\setcounter{figure}{0}
	\setcounter{table}{0}
	\chapter*{\thechapter.\enskip #2}
	\addcontentsline{toc}{chapter}{\thechapter.\enskip #2}
	\markboth{#2}{#2}
}

% Apéndices
\renewcommand{\appendixname}{Apéndice}
\renewcommand*\cftappendixname{\appendixname}

\newcommand{\apendice}[1]{
	%\renewcommand{\thechapter}{A}
	\chapter{#1}
}

\renewcommand*\cftappendixname{\appendixname\ }

% Formato de portada
\makeatletter
\usepackage{xcolor}
\newcommand{\tutor}[1]{\def\@tutor{#1}}
\newcommand{\course}[1]{\def\@course{#1}}
\definecolor{cpardoBox}{HTML}{E6E6FF}
\def\maketitle{
  \null
  \thispagestyle{empty}
  % Cabecera ----------------
\noindent\includegraphics[width=\textwidth]{cabecera}\vspace{1cm}%
  \vfill
  % Título proyecto y escudo informática ----------------
  \colorbox{cpardoBox}{%
    \begin{minipage}{.8\textwidth}
      \vspace{.5cm}\Large
      \begin{center}
      \textbf{TFG del Grado en Ingeniería Informática}\vspace{.6cm}\\
      \textbf{\LARGE\@title{}}
      \end{center}
      \vspace{.2cm}
    \end{minipage}

  }%
  \hfill\begin{minipage}{.20\textwidth}
    \includegraphics[width=\textwidth]{escudoInfor}
  \end{minipage}
  \vfill
  % Datos de alumno, curso y tutores ------------------
  \begin{center}%
  {%
    \noindent\LARGE
    Presentado por \@author{}\\ 
    en Universidad de Burgos --- \@date{}\\
    Tutores: \@tutor{}\\
  }%
  \end{center}%
  \null
  \cleardoublepage
  }
\makeatother

\newcommand{\nombre}{Daniel Drefs Fernandes} %%% cambio de comando

% Datos de portada
\title{Decision Tree Simulator}
\author{\nombre}
\tutor{Carlos López Nozal \\ Ismael Ramos Pérez}
\date{\today}

\begin{document}

\maketitle


\newpage\null\thispagestyle{empty}\newpage


%%%%%%%%%%%%%%%%%%%%%%%%%%%%%%%%%%%%%%%%%%%%%%%%%%%%%%%%%%%%%%%%%%%%%%%%%%%%%%%%%%%%%%%%
\thispagestyle{empty}


\noindent\includegraphics[width=\textwidth]{cabecera}\vspace{1cm}

\noindent D. Carlos López Nozal y Ismael Ramos Pérez, profesores del departamento de Ingeniería Informática, área de Lenguajes y Sistemas Informáticos.

\noindent Exponen:

\noindent Que el alumno D. \nombre, con DNI L1TLL31X4, ha realizado el Trabajo final de Grado en Ingeniería Informática titulado Decision Tree Simulator. 

\noindent Y que dicho trabajo ha sido realizado por el alumno bajo la dirección de los que suscriben, en virtud de lo cual se autoriza su presentación y defensa.

\begin{center} %\large
En Burgos, {\large \today}
\end{center}

\vfill\vfill\vfill

% Author and supervisor
\begin{minipage}{0.45\textwidth}
\begin{flushleft} %\large
Vº. Bº. del Tutor:\\[2cm]
D. nombre tutor
\end{flushleft}
\end{minipage}
\hfill
\begin{minipage}{0.45\textwidth}
\begin{flushleft} %\large
Vº. Bº. del co-tutor:\\[2cm]
D. nombre co-tutor
\end{flushleft}
\end{minipage}
\hfill

\vfill

% para casos con solo un tutor comentar lo anterior
% y descomentar lo siguiente
%Vº. Bº. del Tutor:\\[2cm]
%D. nombre tutor


\newpage\null\thispagestyle{empty}\newpage




\frontmatter

% Abstract en castellano
\renewcommand*\abstractname{Resumen}
\begin{abstract}
\end{abstract}

\renewcommand*\abstractname{Descriptores}
\begin{abstract}
\end{abstract}

\clearpage

% Abstract en inglés
\renewcommand*\abstractname{Abstract}
\begin{abstract}
The aim of this project has been to create a web application that helps in learning the concept of decision trees, their creation and all the prerequisites that are necessary to understand the topic, like entropy and conditional entropy calculation.

The entropy and conditional entropy calculators are driven by user-chosen input values.

The decision tree algorithm that is taught through this application is the ID3 algorithm. A step-by-step visualization is provided to ensure a simple and intuitive learning experience. The visualization includes a dynamically created decision tree, displayed with Scalable Vector Graphics (SVG), and two tables that represent the dataset and relevant values at each step, respectively.

This project has been deployed through the GitHub Pages service at the following URL: \url{https://danieldf01.github.io}.
\end{abstract}

\renewcommand*\abstractname{Keywords}
\begin{abstract}
Decision trees, entropy, conditional entropy, ID3 algorithm.
\end{abstract}

\clearpage

% Indices
\tableofcontents

\clearpage

\listoffigures

\clearpage

\listoftables
\clearpage


\mainmatter
\capitulo{1}{Introduction}

Decision trees belong to the most popular models in machine learning and data mining, serving as intuitive models for decision-making. The algorithm known as ID3, or Iterative Dichotomiser 3, recursively partitions data based on attribute values with the aim of maximizing information gain at each step. It is one of the most popular decision tree algorithms and has significantly contributed to their widespread use.

This project has drawn great inspiration from the ``Seshat Tool''\footnote{Seshat Tool: \url{http://cgosorio.es/Seshat/}} and the research article~\cite{https://doi.org/10.1002/cae.22036} based on it. Its aim is to facilitate users' learning experiences about a concept through intuitive step-by-step simulations of certain algorithms. In the case of Seshat Tool, the target concepts were lexical analysis algorithms. While the main influences are its layout and overall functionality of the presented web application, the related article also carried out a study on its actual effectiveness.

Figure \ref{fig:seshat_student_tests} shows a table out of the article that presents results of a questionnaire that students were asked to do. They took the questionnaire after they ``were taught the main basic concepts of automata and regular expressions, as well as''~\cite{https://doi.org/10.1002/cae.22036} two algorithms that are also taught in the web application. The students were split into two groups, one that would take the test only based on the knowledge gained by a prior theoretical lecture and another that would additionally be granted access to the Seshat Tool.
\imagen{seshat_student_tests}{Results of students' tests}{1}
It can be observed that the group that had access to the tool achieved noticeably better results overall than the other group.

Following the success of the Seshat Tool, this project's aim has been to create a similarly helpful and easy-to-understand decision tree simulator in form of a web application that teaches the ID3 algorithm and all necessary surrounding concepts.
\capitulo{2}{Project objectives}

The primary objective of this project has been to create an interactive and informative web application focused on educating users about the decision tree algorithm ID3.
To achieve this goal, the user would first have to be explained the concepts of entropy and conditional entropy which are prerequisites for understanding decision trees and their creation. 

That is why the first part of the application serves as a tool with which users can learn about entropy and how it is calculated. Alongside cards that hold general information about entropy, the user is provided with a calculator that lets them calculate the entropy for self-chosen values and a visual representation of the binary entropy function.

The conditional entropy calculator works in the same way. Users can input different values that represent the number of instances for an example category of an example attribute and calculate the resulting conditional entropy.

The main focus of the decision tree simulator is the dynamic decision tree creation that makes use of Scalable Vector Graphics (SVG) and whose values are based on the underlying dataset. A data table shows the entirety of the currently loaded dataset. The user can choose between utilizing one of the several example datasets provided by the application and their own datasets that have to be in a CSV file format. The step-by-step visualization of the algorithm provides the user with another table that displays relevant values at each step, e.g., an attribute's information gain. In addition to that, the application presents visual cues that signal the current state of the algorithm at each step.
\capitulo{3}{Theoretical concepts}

In the following, all the theoretical concepts relevant for the understanding of the project will be explained.

\section{Decision Trees}

A decision tree ~\cite{decisiontrees} is a versatile supervised learning algorithm used for classification and regression tasks. Its goal is to predict the value of a variable based on previously processed input. Its hierarchical structure includes a root node and several internal and leaf nodes.

It starts at the root, which represents the feature that best separates the underlying dataset based on a certain criterion. An example would be information gain, which will be explained in a later section. From there branches extend to internal nodes, also called decision nodes. These internal nodes also represent features along with a decision rule that tells us how to further split the data. These features are continuously evaluated until homogenous subsets are created by the leaf nodes. These represent all the possible outcomes of the dataset with each one corresponding to a class label.

Decision tree learning utilizes a divide and conquer approach, iteratively finding the best split points until all or most of the input data is classified.

An example of a decision tree which evaluates whether a person should buy a game or not:
\imagen{decision_tree_example}{Decision tree example}{1}

\subsection{Subsecciones}

Además de secciones tenemos subsecciones.

\subsubsection{Subsubsecciones}

Y subsecciones. 

\section{Entropy} \label{entropy}
In the context of decision tree algorithms, entropy ~\cite{entropy_dash} can be viewed as a measure of impurity in a dataset. It can also be described as a measure of disorder, uncertainty or the expected surprise of a classification, but going forward, the term impurity will be used to avoid confusion. In datasets with binary classes, where variables can only have two possible outcome values, the entropy value lies between 0 and 1, inclusive. The higher the entropy, the more impure the dataset is. In a binary-class dataset, a node that has an equal distribution of, e.g., 5 instances belonging to one class and the other 5 instances belonging to the other class, would have an entropy value of 1. Inversely, a node that has all its instances belong to only one class would have an entropy value of 0, making´it a pure node.
The value is calculated using the following formula:


\[ E(X) = -\sum_{i=1}^{n} {p_i} {log_2 (p_i)} \]


$E(X)$ is the entropy of dataset $X$, $n$ describes the number of classes in the dataset, and $p_i$ the proportion of instances in class $i$ or, in other words, the probability of an instance belonging to class $i$.

This following graph ~\cite{entropy_T} that describes the entropy function in relation to the composition of a node shows really well how entropy is used to determine how impure a node is:
\imagen{entropy_graph}{Descriptive graph of the entropy function}{.7}
In a dataset with binary classes, the x-axis describes the amount of instances in a dataset belonging to the positive class while the y-axis measures their entropy values. It records the lowest values at the extremes where there are only or no positive instances recorded in a given dataset while the highest impurity is reached when the numbers of positive and negative instances are equal.
However, it must be mentioned that the value of entropy can reach values higher than 1 if the dataset possesses more than 2 class labels. 


\subsection{Conditional Entropy}
In the context of decision trees, the conditional entropy~\cite{conditional_entropy_wiki} $E(Y\mid X)$measures the uncertainty in the target variable $Y$, which is usually the class labels, given a certain value of the attribute $X$ that is used for splitting a node.
It can be described as the weighted sum of $E(Y\mid X = x)$ for every possible value $x$ of $X$:
\[E(Y\mid X) = \sum_{x \in X} {p(x)} {E(Y\mid X = x)}\]
The formula sums up the conditional entropies for all the possible values $x$ of attribute $X$. As the weight, it uses $p(x)$, which is the proportion between the number of instances that have the attribute value $x$ and the size of the dataset.
Calculating the conditional entropies for the individual values $x$ of attribute $X$ is possible with the following formula:
\[ E(Y\mid X = x) = -\sum_{y \in Y} {p(y\mid x)} {log_2 {(p(y\mid x))}} \]
$p(y\mid x)$ represents the proportion of instances in the dataset with attribute value $x$ that also have the class value $y$. It sums over all the possible values of the target variable $Y$. It can be used synonymous with the formula for $E(X)$ that was previously shown in \ref{entropy}

In the end, this conditional entropy calculation evaluates the effectiveness of splitting a node based on the particular attribute $X$. The lower the conditional entropy, the lower the impurity of the data after splitting the node based on the particular attribute the conditional entropy was calculated for.


\section{Machine Learning}
Machine learning~\cite{machine_learning} is a field of artificial intelligence which focuses on developing algorithms that are able to learn from given input data and, based on that learned knowledge, generalize to unseen data. Machine learning problems can be categorized into three main tasks: Classification problems, Regression problems, and Clustering problems.

\subsection{Classification problems} \label{classification_problems}
Classification problems~\cite{classification_regression} in machine learning describe the process of predicting a class or category of a given input. The labeled input data is learned with all its attributes and values to reach the goal of building a model that can also categorize new data, that has not been seen before, into one of the given classes or categories. An example would be classifying emails as spam or not spam. This type of machine learning problem will be discussed in more detail in the section \ref{classification}.

\subsection{Regression problems}
Regression problems~\cite{classification_regression} in machine learning describe the process of predicting a continuous value based on labeled input data. Instead of predicting a defined class or category, the goal here is to build a model that can predict the quantity of something. An example would be predicting the salary of a person based on their education degree and previous work experience.

\subsection{Clustering problems}
Clustering problems~\cite{clustering} in machine learning do not aim to predict something, like the previous two. Given some unlabeled input data, the goal is to group similar data points together based on certain features to generate an output of clusters inside of which the data of points are more similar to each other than to data points in other clusters.

\subsection{Decision Trees}
With decision trees~\cite{decisiontrees} ultimately following the goal to create a model that predicts the value of a target variable and doing so by learning labeled input data, they are used for both classification and regression tasks.


\section{Classification} \label{classification}
As mentioned before in \ref{classification_problems}, the goal in classification tasks~\cite{classification_regression} is to assign a class or a label to some input data. Classification aims to create a model that can also generalize to new data. To achieve this, there are multiple algorithms and evaluation metrics that are used to measure a model's performance.

\subsection{Classification algorithms}

\subsubsection{K-Nearest Neighbors}
K-Nearest Neighbors~\cite{knn}, or short KNN, is an algorithm that uses distance as a measure to determine K of a data point's nearest neighbors. Based on that information, a classification or prediction is made about which of the groups the data point is grouped with. KNN does not go through a training stage, instead it just stores and memorizes the training dataset.

\subsubsection{Support Vector Machine}
Support Vector Machine~\cite{svm}, or short SVM, is an algorithm that can perform linear as well as non-linear classification. Linear classifiers assume that a data set is linearly separable while non-linear classifiers are not bound to that restrictive belief. With SVM being able to perform both methods, it is able to find a hyperplane that best separates different classes in a data set by maximizing the margin between classes. 

\subsubsection{Logistic Regression}
This algorithm~\cite{logistic_regression} is mostly used for binary classification problems. It uses the logistic function to determine the probability of a given input belonging to a certain class.

\subsection{Evaluation metrics}

\subsubsection{Confusion Matrix}
A confusion matrix~\cite{evaluation_metrics} is a table that is used to describe a classification model's performance. It presents an overview of the predictions on the test data set against the actual classes:
\imagen{confusion_matrix}{Example of a confusion matrix}{.5}
Where:
TP = True Positive; TN = True Negative; FP = False Positive; FN = False Negative.
It is useful for measuring other metrics like Accuracy, Precision and Recall

\subsubsection{Accuracy}
Accuracy~\cite{evaluation_metrics} measures how many of all the classified instances were given the correct label.
It uses the following formula:
\[ Accuracy = {TP + TN}/{TP + TN + FP + FN} \]
This metric is useful for well balanced data sets. If 99 of the data set's instances belonged to one class and only 1 belonged to the other, likeliness is high that the used model would always predict the class with the higher amount of instances, resulting in a 99\% accuracy due to the data set's composure.

\subsubsection{Precision}
Precision~\cite{evaluation_metrics} describes the ratio of correctly predicted positive cases to the total predicted positive cases. In short words, it measures the accuracy of the positive class predictions and is calculated with the following formula:
\[ Precision = {TP}/{TP + FP} \]
Precision is useful for cases where False Positives are a bigger concern than False Negatives

\subsubsection{Recall (Sensitivity}
Recall~\cite{evaluation_metrics} describes how many of the actual positive cases our model was able to correctly predict. It measures the ability of the model to find all positive instances and is calculated using the following formula:
\[ Recall = {TP}/{TP + FN} \]
Recall is useful when False Negatives are of a higher concern than False Positives.


\section{Referencias}

Las referencias se incluyen en el texto usando cite~\cite{wiki:latex}. Para citar webs, artículos o libros~\cite{koza92}, si se desean citar más de uno en el mismo lugar~\cite{bortolot2005, koza92}.


\section{Imágenes}

Se pueden incluir imágenes con los comandos standard de \LaTeX, pero esta plantilla dispone de comandos propios como por ejemplo el siguiente:

\imagen{escudoInfor}{Autómata para una expresión vacía}{.5}



\section{Listas de items}

Existen tres posibilidades:

\begin{itemize}
	\item primer item.
	\item segundo item.
\end{itemize}

\begin{enumerate}
	\item primer item.
	\item segundo item.
\end{enumerate}

\begin{description}
	\item[Primer item] más información sobre el primer item.
	\item[Segundo item] más información sobre el segundo item.
\end{description}
	
\begin{itemize}
\item 
\end{itemize}

\section{Tablas}

Igualmente se pueden usar los comandos específicos de \LaTeX o bien usar alguno de los comandos de la plantilla.

\tablaSmall{Herramientas y tecnologías utilizadas en cada parte del proyecto}{l c c c c}{herramientasportipodeuso}
{ \multicolumn{1}{l}{Herramientas} & App AngularJS & API REST & BD & Memoria \\}{ 
HTML5 & X & & &\\
CSS3 & X & & &\\
BOOTSTRAP & X & & &\\
JavaScript & X & & &\\
AngularJS & X & & &\\
Bower & X & & &\\
PHP & & X & &\\
Karma + Jasmine & X & & &\\
Slim framework & & X & &\\
Idiorm & & X & &\\
Composer & & X & &\\
JSON & X & X & &\\
PhpStorm & X & X & &\\
MySQL & & & X &\\
PhpMyAdmin & & & X &\\
Git + BitBucket & X & X & X & X\\
Mik\TeX{} & & & & X\\
\TeX{}Maker & & & & X\\
Astah & & & & X\\
Balsamiq Mockups & X & & &\\
VersionOne & X & X & X & X\\
} 
\capitulo{4}{Técnicas y herramientas}

Esta parte de la memoria tiene como objetivo presentar las técnicas metodológicas y las herramientas de desarrollo que se han utilizado para llevar a cabo el proyecto. Si se han estudiado diferentes alternativas de metodologías, herramientas, bibliotecas se puede hacer un resumen de los aspectos más destacados de cada alternativa, incluyendo comparativas entre las distintas opciones y una justificación de las elecciones realizadas. 
No se pretende que este apartado se convierta en un capítulo de un libro dedicado a cada una de las alternativas, sino comentar los aspectos más destacados de cada opción, con un repaso somero a los fundamentos esenciales y referencias bibliográficas para que el lector pueda ampliar su conocimiento sobre el tema.

\section{Bootstrap}
Bootstrap ~\cite{bootstrap5_tutorial} is a front-end framework which is known for providing useful and easy-to-use HTML and CSS templates like tables, buttons, forms and many others. It also comes with JavaScript components like modal dialogues and dropdown menus.

One of Bootstrap's key features is its grid system which lets users divide their web page's contents into rows and columns:
\imagen{bootstrap_grid_system}{The Bootstrap grid system}{.9}
Each row possesses 12 columns which the user can freely utilize to organize the contents that are to be displayed. Bootstrap also automatically puts the device on which the web page is displayed into one of six different size categories, based on the device's screen width. This, in combination with the grid system, allows the developer to make their website responsive to different screens, ranging from large desktop monitors to smartphones.

In addition to Bootstrap's ability to create responsive websites, the wide variety of CSS classes used for common HTML elements like buttons offer further simplicity to the web design aspect of creating a page. With this, maintaining a visual consistency throughout the project is made easier, too.

In this project, the current newest version of Bootstrap, Bootstrap 5, is used. The main advantages are its usage of vanilla JavaScript for its components instead of relying on jQuery, new components for better customization and simplified CSS which reduces file size and loading times for the created pages.
\capitulo{5}{Relevant aspects of the project development}

This chapter focuses on important parts of the project development that had a significant influence on the outcome of the work.

\section{Start of the project}
At the start of the project, there was not much experience in web development. For that reason, the first sprints were not only used to learn about the concepts relevant to the topic of decision trees and start implementing first prototypes, but also to simultaneously gain an understanding of web development using HTML, JavaScript, CSS and surrounding tools. That rhythm continued throughout most of the project as almost each step introduced techniques that were new, but had to be learned and then implemented. This slowed down overall development tempo to some extent.

\section{Decision tree creation without D3.js}
As mentioned in \ref{d3js}, D3.js is an open-source JavaScript library which, among other things, allows the user to dynamically create and change DOM contents based on underlying data. For that reason, at the beginning of development, it was chosen to help create the graph for the Entropy function and later, dynamically create and change the decision trees examples.

However, as development of the project went on and it was time to create a prototype for the creation of decision trees, the plan of using D3.js as the main tool to reach that goal was discarded. This decision was made after an extensive web research on how SVGs displaying trees were created with the library and how they would end up looking. With a few exceptions showing slight deviations, the common structure of D3.js trees was as figure \ref{fig:d3js_tree} displays.
\imagen{d3js_tree}{SVG displaying a tree, created with D3.js}{1}
Since this specific library was not exactly easy to use in the first place, trying to make it look the intended way (displaying labels on the edges and other important values connected to the nodes) would have made the creation significantly harder, if it had even been possible at all.

With that in mind, the decision was made not to use D3.js to create the decision trees, but to assemble the decision trees from simple SVG elements to form templates for branches, decision and leaf nodes, and to develop an algorithm that calculates the tree's elements' relative position to each other in an SVG that is responsive to different screen sizes.

Even though that solution meant more work, it allowed the decision trees to look the intended way in the end, as figure \ref{fig:decision_tree_webapp} shows.
\imagen{decision_tree_webapp}{Dynamically created decision tree using the web application}{.6}

\section{Evaluation of code quality with Codacy before reviews}
As mentioned in \ref{codacy}, Codacy is a code analysis tool that helps developers to improve the quality of their code. With the biweekly meetings having functioned not only as a way to figure out the tasks for the next sprint, but also as an opportunity to review what had been done so far, there was often no time to check the code quality. Codacy provided everything needed to improve the code condition even before the review with the tutors.

It offered the opportunity of checking the code's quality without having to wait for a review meeting or ask the tutors whether a certain way of implementing something was good or not.
This gave the review meetings more space to focus on what had been done rather than how it had been done. Especially during the meetings towards the end of development, this extra time was needed to discuss additional changes that needed to be made before the project was finalized.

\section{Decision between two features}
As the deadline moved closer and the development phase was entering its final phase, it was clear that not all features that were initally set as objectives were going to be implemented. This was due to time constraints brought upon by the workload of other classes and exams.

Two features, however, were still talked about in the final weeks before completion. One was to make the datasets for the decision tree simulator interactive, meaning that the user would be able to dynamically add columns and rows to a dataset that was already loaded. The other was to let the user load own datasets in the form of Comma-separated values (CSV) files to let the application build a decision tree upon.

While it was clear that with the remaining time, it would not be possible to implement both features, a decision had to be made and it was made in favor of the latter. This feature would give the user more freedom by letting them choose a dataset from the internet or create their own and use it to observe how different compositions of datasets would influence the resulting decision tree.

\section{Limitations of a small web application}
With the development nearing its completion and the resulting web application having been a rather small one, it was clear that the program would have some limitations. One being the fact that, even though users are able to load their own datasets in CSV file format, there are some restrictions. E.g., the file cannot contain more than 150 rows of instances. While computational complexity is also part of the reason for this constraint, the main reason for the limit of 150 rows is the resulting decision tree's size. With a tree that is too big, the user might not be able to read individual nodes' contents and therefore have a hard time understanding the algorithm, which would defeat the whole purpose of this project.

Another drawback that the size of the application brings with it is the possibility of branches or their labels overlapping with other nodes. Figure \ref{fig:decision_tree_overlapping} shows how the branch label ``very small'' overlaps with the leaf node ``Leaf 8''.
\imagen{decision_tree_overlapping}{Branch label overlapping with leaf node}{.2}
Due to time constraints, this is a problem that could not be given higher priority for resolution than other necessary refinements, as it would have brought a higher level of workload than was acceptable at the time.
\capitulo{6}{Related works}

Este apartado sería parecido a un estado del arte de una tesis o tesina. En un trabajo final grado no parece obligada su presencia, aunque se puede dejar a juicio del tutor el incluir un pequeño resumen comentado de los trabajos y proyectos ya realizados en el campo del proyecto en curso. 
\capitulo{7}{Conclusions and future lines of work}

\section{Conclusions}
One of the hurdles along the way was to acquire knowledge about the technologies and tools that were necessary to carry out this project in a satisfactory way. This was successfully achieved for technologies like JavaScript, CSS, Bootstrap, D3.js, and others that were relevant for web development.
With that, the main objectives were completed in a way that realizes the initial vision of the project. A web application was built that teaches users the concept of decision trees, their creation through the ID3 algorithm, and surrounding concepts like entropy.
It is a tool that can help students make this topic more tangible and simple to understand during their studies.

\section{Future lines of work}
\subsection{Questionnaires}
Much like in the research article~\cite{https://doi.org/10.1002/cae.22036} about the ``Seshat Tool'' that was mentioned in the introduction, a similar study could be made about the effectiveness of this decision tree simulator. Two groups of students could be asked to take the same questionnaire about the topics of entropy, decision trees and their creation using the ID3 algorithm. One group would be prepared only through theoretical input in form of a lecture, the other would additionally be given access to this application. The results would be analyzed to see if there was a difference in test scores between the two groups.

\subsection{C4.5 algorithm}
The C4.5 algorithm~\cite{c4_5_wiki} is another decision tree algorithm that extends ID3 by being able to handle both continuous and discrete attributes, handle missing attribute values, and prune trees to avoid overfitting. A future work could expand upon this application and implement a simulator for the C4.5 algorithm.

Furthermore, a scatter plot could be used to present how a tree with numerical values works. In this way, the user's understanding of that type of decision trees could be strengthened while additionally, the concept of decision boundaries could be explained in an interactive way.

\subsection{Regression trees}
This project focused solely on decision trees. A future addition to this application could implement a simulator that teaches the concept and creation of regression trees. E.g., the Classification And Regression Tree (CART) algorithm~\cite{cart}, which, as the name suggests, is an algorithm that is applicable for classification and regression problems, could be taught interactively.


\bibliographystyle{plain}
\bibliography{bibliografia}

\end{document}